% GROK NOTE::FL2 independent audit::2026-10-12
\documentclass[11pt]{article}
\usepackage{fontspec}
\usepackage[margin=0.9in]{geometry}
\usepackage{amsmath,amssymb,booktabs,xurl,hyperref}
\title{GROK FL2: Independent Adversarial Audit of FINAL-LINK2}
\author{Grok Bot}
\date{12 October 2026 (PT)}
\begin{document}\maketitle
\paragraph{Status sentence.}
P versus NP remains OPEN: FINAL-LINK2 (fd6aad65, 76d35839) does not close the SAT-to-charged-model compiler arrow; it proves exact Bellman/cycle bookkeeping and a one-scan shared-acquisition theorem for a \emph{given} route trie, but supplies no all-CNF compiler from an arbitrary SAT decider into the charged $(T,S,R)$ model, so the compiler arrow is only partially sharpened and stays OPEN, as does the product-span obstruction.

\section{Scope}
Audited: commit \texttt{fd6aad65} (34 paths: \texttt{main9.tex}, \texttt{final\_link2\_*}, lead/proof/red/research reports) and \texttt{76d35839} (manifest and log only, no mathematical content). A commit search on main for \texttt{FINAL-LINK} after 2026-10-08 returned nothing; the 2026-10-12 CP22LIFT1 commits concern Cutting Planes proof transport and do not revise FINAL-LINK2. All code here is my own (\texttt{/workspace/h\_fl2\_1012/fl2check.py}); no repository script was run.

\section{What FINAL-LINK2 actually claims}
\begin{itemize}
\item PROVED/EXACT (checked by reading): the Bellman path characterization, the cycle-contraction criterion, and the two-wrapper closed form. These are statements \emph{about} a recurrence once a compiler with declared $(R_i,w_e,\theta_e,P_i,h_e)$ exists.
\item PROVED/EXACT: the first-true-literal clause trie is a disjoint, witness-preserving cover; for a \emph{fixed} trie, one source scan builds a shared object $H$ from which every residual is recovered without rereading.
\item CONDITIONAL: the points $(6/5,1/2)$ and $(1,3/4)$ below $\gamma=2\cos(\pi/7)$. The report itself states that no all-CNF construction realizes them.
\item The repository's own status is OPEN. The BatchAcquire report marks ``leaf identifier is the small child'' as falsified as a closure claim.
\end{itemize}

\section{Quantifier and hidden-assumption audit}
\begin{enumerate}
\item \textbf{Main theorem is an implication with an existential premise.} It reads ``\emph{suppose} a fixed uniform compiler gives every CNF an epoch plan with favorable ledgers''. The premise is the open arrow. Proving the implication does not prove the premise.
\item \textbf{Declared, not arbitrary, solver.} The ledgers take a solver exponent $k$ (e.g.\ $k=2$) as given. Nothing turns an arbitrary SAT decider, read as a black box, into charged $(T,S,R)$.
\item \textbf{Route trie is an input.} BatchAcquire ``receives an ordinary CNF and an exact disjunctive restriction cover''. Exactness holds for every trie, but the resource claim depends on a chosen depth $d$. No rule for $d$ makes children contract with a polynomial fanout.
\item \textbf{Promise filter.} Only clauses of width $\le 8$ are candidates. Outside that domain the trie is trivial.
\item \textbf{Interface.} A $(H,u)$ child is not an ordinary CNF, so an ordinary-CNF SAT bound does not apply to it. The report concedes this.
\item \textbf{Validation work is excluded from cost.} That is legitimate as oracle work, but it means the constructor ledger proves nothing about verification cost.
\end{enumerate}

\section{Reproduction (COMPUTATIONAL)}
From my own reimplementation of the greedy selection and the first-true-literal trie, the encoded size $M$, the leaf count $q$, the source reads ($2\times$ literals), and the saving ratio $q/2$ match exactly for PHP$(5,4)$, PHP$(6,5)$ and PHP$(7,6)$: $M=830,1490,2862$; $q=64,125,216$; reads $200,360,588$; save $32,62.5,108$. The ``emitted'' and ``max child'' columns (6008/332, 31990/752, 140952/1653) were \emph{not} reproduced. My plain canonical restriction gives larger values (25708/617, 104885/1219, 381985/2476), so the difference comes from the repository's FINAL-LINK1 \texttt{canonical} routine, which I did not reuse. This affects only constants and works in FINAL-LINK2's favor.

\section{Hostile controls (COMPUTATIONAL)}
\textbf{Fixed depth $d=3$ gives no contraction.} On PHP$(p+1,p)$ the fanout is $q=p^3$ (polynomial), but the scale exponent $\theta=\log|\text{max child}|/\log M$ climbs toward 1: $0.956, 0.973, 0.982, 0.987, 0.990, 0.993$ for $p=4,\dots,9$. At $p=9$, $q=8$ only, because greedy picks an at-most-one clause. Unit-scale steps cannot contract.

\textbf{Halving the child needs $d\approx p-1$, and then the fanout is factorial.} With sound conflict pruning, the smallest depth giving max child $\le M/2$ yields $q=24,120,720,5040,40320$ for $p=4,\dots,8$, which is exactly $p!$. Emitted bits reach $1.19\times10^8$ at PHP$(9,8)$. This is the FINAL-LINK1 superpolynomial pathology, $(p-1)!$ versus $p!$, moved into the fanout term $f_e$. It is not fixed.

\textbf{Sign-flipped planted 3-CNF, ratio 4.26.} For $n=30,60,90$ at $d=3,6,9$, $q=3^d$ and $\theta$ falls only slowly: at $n=90$, $\theta=0.993, 0.985, 0.976$ while $q=27, 729, 19683$. Getting a constant $\theta<1$ requires $d=\Theta(n)$, hence $3^{\Theta(n)}$ fanout.

\textbf{Outside the promise.} A width-9 formula gives $q=1$ and $\theta=1$: no epoch at all.

\section{Verdict}
The compiler arrow is \textbf{OPEN}. FINAL-LINK2 sharpens the target but does not touch its load-bearing quantifier. Its mathematics is correct as far as I checked, and I found no falsification of any theorem it states. I found one error of emphasis: its ``Bellman-sufficient'' framing could be misread as closure, while the cost-relevant step, constructing the epoch plan for every CNF, remains assumed. On pigeonhole and planted controls, contraction and polynomial fanout trade off: they are not achieved together on any tested family. The product-span obstruction is OPEN, and P versus NP remains OPEN.
\end{document}
